\subsection{\problemblock{*Units/Fmt} Data Block}\label{sec:block-units-fmt}

All TAOS output variables have default units and print formats listed in the
Appendix. The \problemblock{*units/fmt} block changes either or both. When a
variable's units are changed, they change everywhere on both input and output. For
example, if altitude is changed from feet to kilometers, every reference to altitude
in table files, problem files, and output files is interpreted in kilometers rather than
feet.

\taossubhead{Units}

Units are changed by giving the variable name and its new unit. Allowed unit tokens
are listed in \cref{tab:allowable-units}. Units must be dimensionally consistent:
a distance variable expressed in feet may be changed to nautical miles, meters, or
kilometers, but not to a unit of time or mass.

\begin{longtable}{@{}>{\ttfamily}lp{0.29\linewidth}>{\ttfamily}lp{0.32\linewidth}@{}}
  \caption{Allowable units.}\label{tab:allowable-units}\\
  \toprule
  \normalfont Units & Description & \normalfont Units & Description \\
  \midrule
  \endfirsthead
  \toprule
  \normalfont Units & Description & \normalfont Units & Description \\
  \midrule
  \endhead
  ft & Feet & m/sec2 & Meters/second$^2$ \\
  in & Inches & m/min2 & Meters/minute$^2$ \\
  mi & Statute miles & m/hr2 & Meters/hour$^2$ \\
  nm & Nautical miles & km/sec2 & Kilometers/second$^2$ \\
  m & Meters & km/min2 & Kilometers/minute$^2$ \\
  km & Kilometers & km/hr2 & Kilometers/hour$^2$ \\
  sec & Seconds & deg/sec2 & Degrees/second$^2$ \\
  min & Minutes & deg/min2 & Degrees/minute$^2$ \\
  hr & Hours & deg/hr2 & Degrees/hour$^2$ \\
  deg & Degrees & rad/sec2 & Radians/second$^2$ \\
  rad & Radians & rad/min2 & Radians/minute$^2$ \\
  ft/sec & Feet/second & rad/hr2 & Radians/hour$^2$ \\
  ft/min & Feet/minute & rev/sec2 & Revolutions/second$^2$ \\
  ft/hr & Feet/hour & rev/min2 & Revolutions/minute$^2$ \\
  in/sec & Inches/second & lb & Pounds mass \\
  in/min & Inches/minute & slugs & Slugs \\
  in/hr & Inches/hour & gm & Grams \\
  mi/sec & Statute miles/second & kg & Kilograms \\
  mi/min & Statute miles/minute & lb/sec & Pounds mass/second \\
  mi/hr & Statute miles/hour & lb/min & Pounds mass/minute \\
  knots & Nautical miles/hour & lb/hr & Pounds mass/hour \\
  m/sec & Meters/second & slugs/sec & Slugs/second \\
  m/min & Meters/minute & slugs/min & Slugs/minute \\
  m/hr & Meters/hour & slugs/hr & Slugs/hour \\
  km/sec & Kilometers/second & g/sec & Grams/second \\
  km/min & Kilometers/minute & g/min & Grams/minute \\
  km/hr & Kilometers/hour & g/hr & Grams/hour \\
  deg/sec & Degrees/second & kg/sec & Kilograms/second \\
  deg/min & Degrees/minute & kg/min & Kilograms/minute \\
  deg/hr & Degrees/hour & kg/hr & Kilograms/hour \\
  rad/sec & Radians/second & lbf & Pounds force \\
  rad/min & Radians/minute & n & Newtons \\
  rad/hr & Radians/hour & kn & Kilonewtons \\
  rev/sec & Revolutions/second & lbf/ft2 & Pounds force/foot$^2$ \\
  rpm & Revolutions/minute & psi & Pounds force/inch$^2$ \\
  g & G's & pascal & Pascals \\
  ft/sec2 & Feet/second$^2$ & kpascal & Kilopascals \\
  ft/min2 & Feet/minute$^2$ & 1/in & Per inch \\
  ft/hr2 & Feet/hour$^2$ & 1/ft & Per foot \\
  in/sec2 & Inches/second$^2$ & 1/m & Per meter \\
  in/min2 & Inches/minute$^2$ & ft2/sec & Feet$^2$/second \\
  in/hr2 & Inches/hour$^2$ & m2/sec & Meters$^2$/second \\
  mi/sec2 & Statute miles/second$^2$ & lb/ft3 & Pounds mass/foot$^3$ \\
  mi/min2 & Statute miles/minute$^2$ & lb/m3 & Pounds mass/meter$^3$ \\
  mi/hr2 & Statute miles/hour$^2$ & g/cm3 & Grams/centimeter$^3$ \\
  nm/hr2 & Nautical miles/hour$^2$ & kg/m3 & Kilograms/meter$^3$ \\
  \bottomrule
\end{longtable}

The aerodynamic reference area \statevar{sref} is not a standard output variable,
but its units can also be changed. Its default is ft$^2$. The distance unit supplied to
the block is squared for area; for example,

\begin{lstlisting}[style=taosmanual]
*units/fmt sref in
\end{lstlisting}

changes the reference-area unit to in$^2$.

\taossubhead{Output Format}

Print format affects output files only, but it applies to every kind of output,
including plot files and EGS databases. A format is written as the letter
\texttt{e} or \texttt{f}, a period, and a number. The letter selects scientific
notation or ordinary fixed-point notation. For example, 250 in E-format is printed as
\texttt{2.50e02}, where \texttt{e02} denotes multiplication by $10^2$; in
F-format it is printed as \texttt{250}. The number after the period gives the number
of digits following the decimal point.

\begin{lstlisting}[style=taosmanual]
*units/fmt alt   km    f.3
           vel   m/sec f.2
           range km
           mach        e.5
\end{lstlisting}

This changes altitude to kilometers with three decimal places, velocity to meters per
second with two decimal places, and range to kilometers with its default format.
Mach number is nondimensional, so only its format is changed, to scientific notation
with five decimal places.

Any number of variables may be changed within one block. Units, format, or both may
be supplied. When both are changed, the unit must appear first and the format second.
The format of a user-defined variable can be controlled in the same way, but TAOS
cannot convert its units because it does not know the variable's original dimensions.
A user-defined variable must therefore be computed in the desired unit system in its
\problemblock{*define} block.
